% Slides 11--14: Siphelele Nkosi, 6:40.

\speakerlabel{Siphelele}{1:40}
\begin{frame}{Heston introduces a second source of risk}
  \begin{columns}[T,onlytextwidth]
    \begin{column}{0.57\textwidth}
      \begin{beamercolorbox}[rounded=true,sep=0.22cm]{block body}
        \[
          \begin{aligned}
          dS_t &= rS_t\,dt+\sqrt{v_t}S_t\,dW_t^S,\\
          dv_t &= \kappa(\theta-v_t)\,dt+\xi\sqrt{v_t}\,dW_t^v,\\
          dW_t^S\,dW_t^v&=\rho\,dt .
          \end{aligned}
        \]
      \end{beamercolorbox}
      \vspace{0.18cm}
      \begin{tightitemize}
        \item \(v_t\): stochastic instantaneous variance.
        \item \(\kappa,\theta\): mean reversion; \(\xi\): volatility of volatility.
        \item \(\rho=-0.70\): correlated spot--variance shocks in the submitted experiment.
        \item \(2\kappa\theta=0.64\geq\xi^2=0.36\): Feller condition satisfied.
      \end{tightitemize}
    \end{column}
    \begin{column}{0.39\textwidth}
      \centering
      \begin{tikzpicture}[
        src/.style={circle,minimum size=0.9cm,fill=Coral!12,draw=Coral!30,font=\small\bfseries},
        asset/.style={rounded corners=4pt,minimum width=2.6cm,minimum height=0.8cm,
          fill=Mist,draw=DeepNavy!20,font=\small\bfseries},
        arr/.style={-{Latex[length=2mm]},line width=0.9pt,color=Teal}
      ]
        \node[src] (ws) at (0,1.55) {\(W^S\)};
        \node[src] (wv) at (2.7,1.55) {\(W^v\)};
        \node[asset] (stock) at (1.35,0.1) {stock};
        \draw[arr] (ws)--(stock);
        \draw[arr,dashed] (wv)--(stock);
        \node[asset,fill=Teal!12] (option) at (1.35,-1.45) {liquid option};
        \draw[arr] (ws)--(option);
        \draw[arr] (wv)--(option);
      \end{tikzpicture}
    \end{column}
  \end{columns}
  \sourcefoot{Submitted report, Secs.\ 2.1.2 and 4.1.4; Heston (1993).}
  \note{Explain the two diffusion risks and why stock plus cash cannot span both. The liquid option is useful because it has variance sensitivity.}
\end{frame}

\speakerlabel{Siphelele}{1:45}
\begin{frame}{The hedge now chooses two traded positions}
  \begin{columns}[T,onlytextwidth]
    \begin{column}{0.52\textwidth}
      {\fontsize{8.8}{10.2}\selectfont\[
        G_T=
        \sum_{n=0}^{N-1}\delta_n\,\Delta S_n
        +\sum_{n=0}^{N-1}\eta_n\,\Delta C^h_n
      \]}
      \vspace{0.06cm}
      \begin{beamercolorbox}[rounded=true,sep=0.18cm]{block body}
        {\fontsize{8.1}{9.4}\selectfont\[
          x_n=\left(
          \log\frac{S_n}{K},\,
          \frac{T-t_n}{T},\,
          \frac{v_n}{\theta},\,
          \log\frac{S_n}{K_h},\,
          \frac{T_h-t_n}{T_h},\,
          \frac{C_n^h}{S_n}
          \right)
        \]}
      \end{beamercolorbox}
      \vspace{0.05cm}
      {\scriptsize Target \(K=0.9,T=0.5\); liquid option \(K_h=1,T_h=1\), alive throughout.}
    \end{column}
    \begin{column}{0.44\textwidth}
      \begin{tikzpicture}[
        box/.style={rounded corners=4pt,fill=Mist,draw=DeepNavy!18,
          minimum width=3.3cm,minimum height=0.62cm,align=center,font=\scriptsize},
        arr/.style={-{Latex[length=2mm]},color=Teal,line width=0.9pt},
        node distance=0.12cm
      ]
        \node[box] (state) {six causal state features};
        \node[box,below=of state] (mlp) {shared fully connected network};
        \node[box,below=of mlp,fill=Teal!12] (out) {\(\delta_n=5\tanh z_{1,n}\)\\\(\eta_n=5\tanh z_{2,n}\)};
        \node[box,below=of out] (gain) {two-instrument gains \(\rightarrow\) terminal MSE};
        \draw[arr] (state)--(mlp);\draw[arr] (mlp)--(out);\draw[arr] (out)--(gain);
      \end{tikzpicture}
      \vspace{0.06cm}
      {\scriptsize\begin{tightitemize}
        \item Signed positions are now economically necessary.
        \item Full-truncation protects variance; all positions are causal.
      \end{tightitemize}}
    \end{column}
  \end{columns}
  \sourcefoot{Submitted report, Sec.\ 4.1.4 and Algorithms A.3.3, A.6.10.}
  \note{Define delta and eta. Explain why sigmoid is no longer appropriate. Mention log-Euler for stock and full-truncation Euler for variance.}
\end{frame}

\speakerlabel{Siphelele}{1:45}
\begin{frame}{The frozen-volatility proxy changed the market}
  \begin{columns}[T,onlytextwidth]
    \begin{column}{0.54\textwidth}
      \vspace{0.05cm}
      \begin{beamercolorbox}[rounded=true,sep=0.28cm]{block body}
        \[
          C_t^h \approx
          C_{\bs}\!\left(S_t,K_h,T_h-t,\sqrt{v_t}\right)
        \]
      \end{beamercolorbox}
      \vspace{0.12cm}
      \claimbox{Insert the current Heston variance into a familiar Black--Scholes price.}
    \end{column}
    \begin{column}{0.42\textwidth}
      \pill{Why it seemed reasonable}
      {\small\begin{tightitemize}
        \item cheap to evaluate at every path state;
        \item responsive to the current variance;
        \item produces a volatility-sensitive second instrument;
        \item preliminary neural RMSE \(0.008198\) was encouraging.
      \end{tightitemize}}
      \vspace{0.05cm}
      \pill[Coral]{The hidden mismatch}
      {\small\begin{tightitemize}
        \item \(S_t,v_t\) follow Heston;
        \item \(C_t^h\) solves a frozen-volatility Black--Scholes equation.
      \end{tightitemize}}
    \end{column}
  \end{columns}
  \vspace{0.06cm}
  \risknote{The proxy experiment motivates the refinement; it is not compared directly with the final COS experiment because the traded-option processes differ.}
  \sourcefoot{Submitted report, Sec.\ 4.1.4 (pp.\ 24--25) and Algorithm A.6.5.}
  \note{The preliminary 0.008198 value is report-supported, but the separate run archive is absent. Keep it secondary. The core point is the model mismatch.}
\end{frame}

\speakerlabel{Siphelele}{1:30}
\begin{frame}{A martingale diagnostic exposes inconsistency}
  \begin{columns}[T,onlytextwidth]
    \begin{column}{0.63\textwidth}
      \begin{tikzpicture}[x=1cm,y=1cm]
        \draw[-{Latex[length=2mm]},Slate] (1.2,0)--(6.35,0);
        \draw[dashed,Slate] (6.0,-0.15)--(6.0,1.85);
        \node[below,font=\scriptsize,color=Slate] at (6.0,-0.1) {zero};
        \node[right,font=\small] at (0,1.35) {BS proxy};
        \node[right,font=\small] at (0,0.45) {Heston COS};
        \draw[line width=10pt,color=Coral!70] (1.35,1.35)--(6.0,1.35);
        \draw[line width=10pt,color=Teal!75] (5.94,0.45)--(6.0,0.45);
        \node[anchor=west,font=\small\bfseries,color=Coral] at (1.35,1.65) {\(-0.009482\)};
        \node[anchor=east,font=\small\bfseries,color=Teal] at (5.9,0.75) {\(-0.000118\)};
        \node[anchor=west,font=\scriptsize,color=Slate] at (1.35,1.08) {SE \(0.000459\)};
        \node[anchor=east,font=\scriptsize,color=Slate] at (5.9,0.18) {SE \(0.000479\)};
      \end{tikzpicture}
    \end{column}
    \begin{column}{0.33\textwidth}
      \begin{tightitemize}
        \item Under \(r=0\), a consistently priced traded option should have approximately zero mean movement.
        \item The proxy displacement is large relative to its standard error.
        \item The COS residual is statistically negligible and consistent with discretisation.
      \end{tightitemize}
    \end{column}
  \end{columns}
  \vspace{0.2cm}
  \claimbox{The proxy is not merely imprecise in level; its path dynamics are inconsistent with the simulated Heston market.}
  \sourcefoot{Submitted report, Table A.9 and Sec.\ 4.1.4. Bars reproduce reported values; zero is the risk-neutral \(r=0\) diagnostic reference.}
  \note{Avoid the stronger phrase ``formal arbitrage proof.'' Handover: ``The diagnostic is empirical; Glasson will show the exact missing drift terms and the model-consistent correction.''}
\end{frame}
